\section{The exact physical Gram expansion}
\label{sec:tpc34-exact-gram}

We first open the two structures that are invisible in the compact matrix
\(K^L\): the ultra divisor \(u\) and the row-gcd projector divisor \(v\).
All sums are finite, so the rearrangements below are identities.

\subsection{A synthesis matrix for the left polarization}

Let \(\mathfrak X_L\) be the set of quadruples
\begin{equation}
 \xi=(\alpha,j,u,v)
 \label{eq:tpc34-xi}
\end{equation}
on the physical support such that
\begin{equation}
 T<u\le U_0,\qquad u\mid N_\alpha(j),\qquad v\mid d_\alpha.
 \label{eq:tpc34-xi-conditions}
\end{equation}
Define
\begin{align}
 z_L(\xi)
 &:=
 \gamma_\alpha^{(1)}a_R(u)\lambda_G(v),
 \label{eq:tpc34-z-left}\\
 \Phi_\gamma^L(\xi)
 &:=
 \one_{v\mid d_\gamma}
 \mathfrak B_{\alpha,\gamma}(j)
 \mathsf H_{m_\gamma}(j).
 \label{eq:tpc34-Phi-left}
\end{align}
The conditions in \eqref{eq:tpc34-xi-conditions} are part of the index
set and therefore are not repeated as indicators in \(z_L\).

\begin{theorem}[Exact opened \(TT^*\) identity]
\label{thm:tpc34-exact-Gram}
Let \(T_L:\C^{\mathfrak X_L}\to\C^{\{\gamma\}}\) be the synthesis matrix
\begin{equation}
 (T_Lz)(\gamma)=\sum_{\xi\in\mathfrak X_L}
 \Phi_\gamma^L(\xi)z(\xi).
 \label{eq:tpc34-synthesis-map}
\end{equation}
Then
\begin{align}
 B_\gamma^L
 &=\sum_{\xi\in\mathfrak X_L}
 \Phi_\gamma^L(\xi)z_L(\xi),
 \label{eq:tpc34-opened-left-column}\\
 E_L
 &=\|T_Lz_L\|_2^2
 \notag\\
 &=\sum_{\xi_1,\xi_2\in\mathfrak X_L}
 \overline{z_L(\xi_1)}z_L(\xi_2)
 \cG_L(\xi_1,\xi_2),
 \label{eq:tpc34-exact-left-Gram}\\
 \cG_L(\xi_1,\xi_2)
 &:=
 \sum_\gamma
 \overline{\Phi_\gamma^L(\xi_1)}
 \Phi_\gamma^L(\xi_2)
 =(T_L^*T_L)(\xi_1,\xi_2).
 \label{eq:tpc34-left-Gram-kernel}
\end{align}
In particular \(\cG_L\) is positive semidefinite.  There is an identical
right-polarized identity after exchanging \(\alpha\) and \(\gamma\).
\end{theorem}

\begin{proof}
Insert \eqref{eq:tpc34-ultra-increment} and
\eqref{eq:tpc34-gcd-projector} into
\eqref{eq:tpc34-left-energy}.  The part of the physical multiplier not
belonging to the row-gcd indicator is exactly
\(\mathfrak B_{\alpha,\gamma}(j)\).  This gives
\eqref{eq:tpc34-opened-left-column}.  Squaring, summing over \(\gamma\),
and interchanging finite sums gives
\eqref{eq:tpc34-exact-left-Gram}.  The last identity is the definition of
the adjoint product.
\end{proof}

\begin{remark}[Which diagonal is meant]
\label{rem:tpc34-two-diagonals}
The actual mask vanishes on the physical pair diagonal \(\alpha=\gamma\).
The new Gram expansion contains two independent copies
\(\alpha_1,\alpha_2\) of the \emph{input} row.  Its diagonal
\(\alpha_1=\alpha_2\) is not deleted by the mask; each copy may still be
different from the common opposite row \(\gamma\).  We call this the
Gram-copy or same-input-row diagonal.
\end{remark}

\subsection{Divisor incidence in the Gram kernel}

\begin{proposition}[Least-common-multiple incidence]
\label{prop:tpc34-Gram-incidence}
Write \(\xi_i=(\alpha_i,j_i,u_i,v_i)\).  Uniformly on the physical packet,
\begin{equation}
 \boxed{
 |\cG_L(\xi_1,\xi_2)|
 \ll_{\eps,\mathscr D}
 X^\eps\left(L+\frac{Q}{[v_1,v_2]}\right).}
 \label{eq:tpc34-Gram-incidence}
\end{equation}
The same estimate holds for the right Gram kernel.
\end{proposition}

\begin{proof}
A nonzero summand in \eqref{eq:tpc34-left-Gram-kernel} requires
\[
 v_1\mid d_\gamma,\qquad v_2\mid d_\gamma,
\]
and hence \([v_1,v_2]\mid d_\gamma\).  For each permitted prime
\(\ell_\gamma\asymp L\), the number of integers
\(d_\gamma\asymp D\) divisible by \([v_1,v_2]\) is at most
\[
 O_{\mathscr D}\!\left(1+\frac{D}{[v_1,v_2]}\right).
\]
There are \(O(L)\) possible prime labels, and the remaining factors in
\(\Phi_\gamma^L\) have an \(X^\eps\) pointwise envelope by
\eqref{eq:tpc34-prefix-envelope}.  Since \(LD=Q\), this proves
\eqref{eq:tpc34-Gram-incidence}.  The proof for the other polarization is
the same.
\end{proof}

\begin{remark}[The principal projector layer]
\label{rem:tpc34-principal-v}
When \(v_1=v_2=1\), the right-hand side of
\eqref{eq:tpc34-Gram-incidence} is \(O(X^\eps Q)\), the full opposite-row
count.  Thus the projector supplies a real gain on large-lcm layers but
does not suppress its principal layer.  The columnwise coefficient-mass
bound \eqref{eq:tpc34-gcd-projector-mass} cannot change this conclusion.
\end{remark}

\subsection{The absolute baseline}

\begin{proposition}[Sign-blind baseline]
\label{prop:tpc34-absolute-baseline}
Without using cancellation between input rows or orbit points,
\begin{equation}
 \boxed{
 E_L+E_R\ll_{\eps,\mathscr D}X^\eps Q^3J^2.}
 \label{eq:tpc34-absolute-baseline}
\end{equation}
\end{proposition}

\begin{proof}
By \eqref{eq:tpc34-prefix-envelope},
\eqref{eq:tpc34-row-l1}, and the \(O(J)\) orbit support,
\[
 |B_\gamma^L|
 \ll X^\eps J\sum_\alpha|\gamma_\alpha^{(1)}|
 \ll X^\eps QJ.
\]
Summing its square over \(O(Q)\) opposite rows proves the left estimate.
The right estimate is symmetric.
\end{proof}

The factor \(J^2\) in \eqref{eq:tpc34-absolute-baseline} comes from
coherent duplication of the two orbit times.  The next section proves
that several geometrically defined subfamilies are much smaller, but also
shows where the unsigned incidence estimates proved here stop.
